% ---------------------------------------------------------------------------
% Chapter 10: The Second Checkpoint: High-Level Trigger
% Generated from the outline; edit freely, this file is now the source.
% ---------------------------------------------------------------------------
\chapter{The Second Checkpoint: High-Level Trigger}
\chaptermark{The Second Checkpoint}
\label{ch:second_checkpoint}

\begin{journeybox}
The second checkpoint is staffed by the same kind of officers who will later
conduct the census, but in a hurry. They reconstruct the family quickly,
calibrate it with an older table, and keep only one crossing in a hundred.
\end{journeybox}

\physicsline{HLT paths, Calo jets, PF jets, PUPPI at HLT (Run~3), online JEC, prescales, trigger efficiency (tag-and-probe, orthogonal datasets)}

\section{From 100~kHz to 1~kHz: HLT design}
\label{sec:second_checkpoint:from_100_khz_to_1_khz_hlt_design}
\todo{Write this section.}

\section{Calo-jet and PF-jet reconstruction online; PUPPI online in Run~3}
\label{sec:second_checkpoint:calo_jet_and_pf_jet_reconstruction_onlin}
\todo{Write this section.}

\section{Online jet energy corrections and why they differ from offline}
\label{sec:second_checkpoint:online_jet_energy_corrections_and_why_th}
\todo{Write this section.}

\section{Measuring trigger efficiency for a jet analysis}
\label{sec:second_checkpoint:measuring_trigger_efficiency_for_a_jet_a}
\todo{Write this section.}

\section{Our jet at this stage: HLT objects and their $\pt$}
\label{sec:second_checkpoint:our_jet_at_this_stage_hlt_objects_and_th}
\todo{Numbers, plots and event display of the $b$-jet at this stage.}

\begin{ledger}{The Second Checkpoint}
  \ledgerrow{before this stage}{}{}{}{--}
  \ledgerrow{after this stage}{}{}{}{}
\end{ledger}

\section{Further reading}
\label{sec:second_checkpoint:further_reading}
Official and theory references for this stage: \cite{CMS:2016ngn}, \cite{CMS:2024aqx}.

\begin{pitfalls}
\todo{Pitfalls and FAQs for newcomers at this stage.}
\end{pitfalls}

\begin{exercises}
\item \todo{Exercise 1.}
\item \todo{Exercise 2.}
\end{exercises}
